\chapter{Impacto, ética e identidad nacional}
\label{cap:impacto}
\section{Impacto tecnológico}
ExploraNica aplica una arquitectura de servicios centralizados a la gestión operativa del turismo. Los clientes de escritorio comparten una API y una base de datos con contratos, estados y relaciones definidos. Esta organización concentra las operaciones sobre una fuente común y permite ampliar los servicios sin duplicar la persistencia en cada cliente.

El proyecto integra programación orientada a objetos, persistencia relacional, autenticación, consumo de servicios, programación asíncrona y control de versiones. La combinación de estas tecnologías permite desarrollar interfaces de gestión y coordinar operaciones sobre datos compartidos. La separación de componentes facilita mantener el producto y ampliar sus servicios.

La sostenibilidad tecnológica se apoya en dependencias mantenidas, decisiones de arquitectura documentadas y separación entre clientes y servidor. Los contratos permiten evolucionar los componentes de presentación y mantener las reglas de negocio en los servicios de la API.

\section{Impacto económico}
ExploraNica organiza la oferta de las agencias y facilita su incorporación a procesos digitales. La Cuenta Satélite de Turismo del Banco Central de Nicaragua describe la producción, el consumo y el aporte económico del sector y sirve como referencia para analizar su contexto \cite{ref2}. El proyecto se orienta a las necesidades de gestión de los prestadores de servicios turísticos.

El modelo de negocio plantea una comisión del 8\,\% por reserva. Para una base de 100 unidades monetarias, esta regla corresponde a una comisión de 8 y un saldo preliminar de 92, antes de costos de pasarela, impuestos y devoluciones. Su aplicación requiere incorporar el módulo financiero, registrar las transacciones y definir las reglas de conciliación.

La evaluación económica futura debe registrar costo de infraestructura, mantenimiento, capacitación y soporte; reservas atendidas; ingresos efectivamente conciliados y tasa de cancelación. Sólo con estos datos podrá valorarse si la herramienta aporta ahorro operativo o rentabilidad suficiente para sostener su servicio.

\section{Impacto social}
La digitalización facilita que las agencias de menor tamaño organicen su oferta con una estructura común de destinos, precios, fechas y responsables. La administración de expedientes y estados establece un procedimiento de revisión compartido. La claridad de la información y la identificación del operador contribuyen a una relación más ordenada entre prestadores y usuarios.

El acceso también puede verse limitado por conectividad, disponibilidad de equipos Windows y experiencia digital. El piloto debe incluir capacitación breve, tareas comprensibles y mecanismos de soporte. La inclusión requiere probar legibilidad, navegación por teclado, contraste y claridad de los mensajes, además de considerar agencias fuera de los centros urbanos.

El sistema complementa la atención humana mediante la organización del servicio y el registro de las decisiones. La atención de incidentes de viaje, disputas y necesidades de accesibilidad permanece a cargo de las personas responsables. La capacitación y el soporte facilitan su incorporación a la operación de las agencias.

\section{Ética profesional en el desarrollo tecnológico}
La ética profesional orienta el diseño de reglas transparentes, la protección de datos y el tratamiento responsable de la información turística. El sistema diferencia las responsabilidades del administrador, el operador y el turista. Las operaciones críticas utilizan controles de identidad y estado para proteger la información y mantener la integridad del servicio.

La protección de datos comprende acceso restringido a los expedientes, uso de información necesaria para prestar el servicio y gestión de su conservación. Las interfaces de consulta limitan los datos documentales expuestos. Las credenciales, los tokens y las claves de firma se administran como secretos de configuración y se excluyen del material de capacitación.

La autorización debe aplicarse en el servidor y comprobar la relación entre usuario y recurso. El equipo debe revisar los accesos por rol y documentar incidentes sin exponer datos personales. La bitácora debe identificar al responsable real, proteger su integridad y evitar almacenar secretos. Los registros actuales con responsable genérico requieren mejorarse para ofrecer trazabilidad individual.

La verificación de agencias comunica su estado y el motivo de rechazo para permitir la corrección del expediente. La incorporación de reseñas y pagos requiere reglas de moderación, cargos, cancelación y reclamación comprensibles para el usuario. La operación del servicio debe integrar los requisitos normativos aplicables al tratamiento de datos y a la actividad turística.

\section{Historia e identidad nacional en el proyecto}
La identidad nacional se expresa en la selección de Nicaragua como territorio de referencia y en la vinculación de sus destinos con operadores locales. Ometepe, Somoto y San Juan del Sur representan escenarios diversos de paisaje y actividad turística. El catálogo permite organizar experiencias asociadas con destinos y prestadores del país.

El catálogo debe presentar la historia y cultura de cada destino con fuentes identificadas y revisión local. Las descripciones de patrimonio, acontecimientos y tradiciones deben evitar errores, estereotipos o apropiación de expresiones comunitarias. Cuando se incorporen fotografías, relatos o conocimientos de una comunidad, debe verificarse su autorización de uso y reconocerse su procedencia.

La utilización del español en los mensajes y de una identidad visual vinculada al país contribuye a la familiaridad del producto. No obstante, los símbolos y nombres nacionales por sí solos no garantizan pertinencia cultural. Esta debe evaluarse mediante participación de operadores y comunidades, respeto por sus formas de representación y actualización responsable de la información.

La proyección turística debe equilibrar promoción y cuidado del entorno: describir recomendaciones de visita, capacidad de atención y condiciones del servicio, sin prometer acceso irrestricto a espacios sensibles. El sistema puede incorporar campos de procedencia y fecha de revisión del contenido para mantener responsabilidad editorial.

\section{Proyección}
La proyección inmediata es consolidar un prototipo demostrable para agencias y administración, con instalación reproducible, pruebas del núcleo y evidencias de uso. La siguiente etapa es un piloto acotado en el que se midan tareas concretas: crear una oferta, programar una salida, consultar reservas y validar atención. Cada participante debe conocer el propósito de la evaluación y el tratamiento de sus datos.

Si el piloto acredita utilidad, puede avanzarse hacia pagos de prueba, reportes y acceso móvil, sin comprometer de antemano una operación comercial a escala nacional. La expansión requiere responsables de soporte, presupuesto de mantenimiento, respaldo y recuperación, acuerdos con operadores y políticas de atención al usuario.

La proyección de largo plazo es consolidar una plataforma de gestión turística con participación local y mejora continua. La expansión del producto se apoya en estabilidad técnica, uso sostenido y sostenibilidad económica. Los nuevos módulos deben responder a necesidades concretas de operadores y turistas, con criterios de aceptación definidos.
